\section{生产力即产品：数据飞轮与真实交付}

上一章说明公司要避开叙事和经济性两个泥潭，本章进一步回答：如果不靠故事，具身智能到底卖什么。大模型可以说“智能即产品”，因为用户直接购买聊天、写作、代码、搜索和 Agent 能力；具身智能则必须说“生产力即产品”。机器人不只是展示智能，而是要在物理世界创造可计量的劳动替代、效率提升或新增产能。客户最终购买的不是“像人”，而是货架整理、搬运、拣选、巡检、清洁、配送等任务的稳定产出。

这也解释了为什么真实数据采集非常贵。雇人遥操作机器人采集真实数据，成本高、速度慢、场景有限；节目描述里提到真实数据在训练数据中只占 1\%，合成数据管线挑起大梁。这个比例不能机械照搬到所有公司，但它说明：如果没有合成数据、仿真和场景内泛化，真实数据成本很难支撑具身智能早期规模化。

\begin{figure}[H]
\centering
\includegraphics[width=0.96\textwidth]{productivity-as-product.png}
\caption{生产力即产品：大模型卖智能，具身智能必须卖可衡量生产力。自制概念图，依据 01:55:17--02:13:51 对谈内容与视频描述整理。}
\end{figure}

\begin{knowledgebox}{读图：从任务到收入必须经过成本}
图中 Task 到 Robot 只是开始，真正关键是 Cost、Output 和 Revenue。机器人要产生产品价值，必须让单位任务成本下降，输出稳定可验，客户愿意复购。没有成本下降的智能，只是昂贵演示。
\end{knowledgebox}

\subsection{遥操作经济账}

本节把“真实数据很贵”具体算出来。王鹤给出的估算是：一台 full-size 人形机器人制造成本至少十万元级；如果为了采数据买一万台，硬件投入就是十亿元级。更贵的是运营：每台机器人如果两班倒，甚至每班两个人遥操作，再加标注、质检和维护，一个月维持万台采集系统的成本会到数亿甚至十亿元量级。这个账解释了为什么“全部靠真实遥操作数据”很难成为早期通用路线。

\begin{knowledgebox}{经济账：为什么不能只靠遥操作采数据}
\begin{tabularx}{\textwidth}{p{0.22\textwidth}p{0.28\textwidth}X}
\toprule
成本项 & 粗略量级 & 启发 \\
\midrule
机器人硬件 & 万台 $\times$ 十万元级 & 采集前就需要十亿元级资产投入。 \\
遥操作人力 & 每台多班次、多人协作 & 人力不是一次性成本，而是持续 burn。 \\
标注与质检 & 轨迹、状态、失败、任务结果都要检查 & 动作数据不能像普通图片一样粗标。 \\
场景与维护 & 场地、维修、补货、异常处理 & 物理世界会不断制造非模型问题。 \\
\bottomrule
\end{tabularx}
\end{knowledgebox}

这也是自动驾驶与具身智能的重要差别。汽车可以卖给用户，用户每天开车时顺带产生数据，数据成本被产品使用摊掉；机器人如果还不能独立创造价值，就无法把数据采集伪装成商业化。把“没有功能的机器人”卖给客户，让客户自己采、自己训，短期也许能形成收入，但长期会把研发风险转嫁给客户。

\subsection{数据 recipe 与出货回流}

本节把数据问题转成飞轮问题。具身智能的数据 recipe 可以理解为：互联网/视觉数据提供常识和表征，合成数据提供长尾和可控训练，少量真实数据提供 grounding，出货后的真实失败提供持续回流。早期没有出货量时，真实数据比例可能很低；一旦进入真实场景，数据回流会把产品和训练管线连起来。

\begin{figure}[H]
\centering
\includegraphics[width=0.92\textwidth]{data-recipe-robotics.png}
\caption{具身数据 Recipe：真实 1\%、合成管线、出货回流和数据飞轮。自制概念图，依据 01:55:17--02:13:51 对谈内容与视频描述整理。}
\end{figure}

\begin{knowledgebox}{读图：越靠上越贴近真实，也越稀缺}
底层通用数据建立基础能力，中层合成数据放大任务分布，少量真实数据校准，出货回流提供最关键的失败和边界案例。读图时要避免一个误解：真实数据占比低不代表不重要，恰恰因为它稀缺，所以要用来校准和验证整个系统。
\end{knowledgebox}

\begin{warningbox}{没有功能的机器人不能卖成数据飞轮}
访谈描述提到一种 tricky 现象：把没有功能的机器人卖给别人，让客户自己采、自己训，并承诺未来会成为员工。这个模式风险很大，因为它把研发风险转嫁给客户，也会透支行业信用。
\end{warningbox}

\begin{importantbox}{生产力的定义}
王鹤把生产力定义得很朴素：单位时间内干的活必须跟人相当，或者至少在一个清晰边界内提升效率。如果机器人比人慢太多、干不久、故障多，或者必须靠大量人工兜底，它就不是优质生产力，甚至可能是落后生产力。
\end{importantbox}

\subsection{本章小结}

具身智能的产品不是“看起来聪明”，而是可计量生产力。数据飞轮也必须建立在真实交付上：合成数据降低早期成本，出货回流提供真实校准，客户价值验证路线是否成立。

